\documentclass[aspectratio=169, 10pt]{beamer}

% ==============================================================================
% GÓI HỖ TRỢ TIẾNG VIỆT, PHÔNG CHỮ & ĐỒ HỌA
% ==============================================================================
\usepackage[utf8]{inputenc}
\usepackage[vietnamese]{babel}
\usepackage{lmodern}
\usepackage{tikz}
\usepackage{tcolorbox}
\usepackage{booktabs}
\usepackage{tabularx}
\usepackage{amsmath}

% ==============================================================================
% BẢNG MÀU DỊU MẮT (GENTLE ACADEMIC PALETTE - XANH NAVY & XANH NGỌC)
% ==============================================================================
\definecolor{ThemeNavy}{RGB}{30, 58, 138}     % Xanh Navy đậm thanh lịch (#1E3A8A)
\definecolor{PrimaryBlue}{RGB}{37, 99, 235}    % Xanh công nghệ (#2563EB)
\definecolor{SoftTeal}{RGB}{13, 148, 136}      % Xanh ngọc dịu mát (#0D9488)
\definecolor{SoftSage}{RGB}{16, 185, 129}      % Xanh lá nhã nhặn (#10B981)
\definecolor{WarmAmber}{RGB}{217, 119, 6}      % Hổ phách cảnh báo nhẹ (#D97706)
\definecolor{MutedSlate}{RGB}{71, 85, 105}     % Xanh xám phiến (#475569)
\definecolor{CardBg}{RGB}{248, 250, 252}       % Nền thẻ sáng dịu (#F8FAFC)
\definecolor{TextDark}{RGB}{15, 23, 42}        % Chữ than chì dễ đọc (#0F172A)

% Cấu hình Theme Beamer
\usetheme{Madrid}
\usecolortheme{default}

\setbeamercolor{palette primary}{bg=ThemeNavy, fg=white}
\setbeamercolor{palette secondary}{bg=SoftTeal, fg=white}
\setbeamercolor{palette tertiary}{bg=PrimaryBlue, fg=white}
\setbeamercolor{structure}{fg=ThemeNavy}
\setbeamercolor{title}{bg=ThemeNavy, fg=white}
\setbeamercolor{frametitle}{bg=ThemeNavy, fg=white}
\setbeamercolor{item}{fg=SoftTeal}
\setbeamercolor{normal text}{fg=TextDark}

% Cấu hình khối hộp tcolorbox phẳng (Flat & Modern)
\tcbset{
    boxrule=0.7pt,
    arc=1.5mm,
    fonttitle=\bfseries\small,
    coltitle=white,
    left=4pt, right=4pt, top=4pt, bottom=4pt
}

% ==============================================================================
% THÔNG TIN TIÊU ĐỀ
% ==============================================================================
\title[IT3120 - User Problem \& Lean UX Canvas]{XÁC ĐỊNH VẤN ĐỀ NGƯỜI DÙNG \& XÂY DỰNG LEAN UX CANVAS}
\subtitle{Nghiên cứu Trải nghiệm Lập Kế hoạch \& Đăng ký Học phần Đại học}
\author[Nhóm Sinh viên]{Nhóm Sinh viên Thực hiện}
\institute[SOICT - HUST]{
    Trường Công nghệ Thông tin \& Truyền thông (SOICT)\\
    Đại học Bách khoa Hà Nội
}
\date{\today}

\begin{document}

% ------------------------------------------------------------------------------
% SLIDE 1: TRANG TIÊU ĐỀ
% ------------------------------------------------------------------------------
\begin{frame}
    \titlepage
\end{frame}

% ------------------------------------------------------------------------------
% SLIDE 2: BỐI CẢNH & 5 MỤC TIÊU CỐT LÕI
% ------------------------------------------------------------------------------
\begin{frame}{1. Bối cảnh Nghiên cứu \& 5 Mục tiêu Bài thực hành}
    \begin{columns}[T]
        \begin{column}{0.48\textwidth}
            \begin{tcolorbox}[colback=CardBg, colframe=ThemeNavy, title=Bối cảnh Nghiên cứu]
                \small
                \textbf{Đại học quy mô lớn (35.000+ sinh viên):}
                \begin{itemize}
                    \item Đào tạo theo học chế tín chỉ: Sinh viên chủ động xây dựng kế hoạch học tập cá nhân.
                    \item \textbf{Khó khăn hiện hữu:} Thông tin môn học phân tán, áp lực nghẽn mạng giờ cao điểm mở cổng.
                    \item \textbf{Hệ quả:} Lịch học bị xáo trộn, gia tăng khối lượng xử lý đơn từ thủ công cho nhà trường.
                \end{itemize}
            \end{tcolorbox}
        \end{column}
        
        \begin{column}{0.48\textwidth}
            \begin{tcolorbox}[colback=CardBg, colframe=SoftTeal, title=5 Mục tiêu Cốt lõi]
                \small
                \begin{enumerate}
                    \item \textbf{Phân biệt rành mạch:} User Needs, User Problem và Business Problem.
                    \item \textbf{Viết Problem Statement} dựa trên bằng chứng nghiên cứu thực tế.
                    \item \textbf{Sử dụng Canvas 6 ô} làm rõ bối cảnh, rào cản và giải pháp thay thế.
                    \item \textbf{Kết nối nhân quả} hai loại vấn đề (không đánh đồng).
                    \item \textbf{Xây dựng Lean UX Canvas} với giả thuyết và kế hoạch kiểm chứng.
                \end{enumerate}
            \end{tcolorbox}
        \end{column}
    \end{columns}
\end{frame}

% ------------------------------------------------------------------------------
% SLIDE 3: PHÂN BIỆT 3 KHÁI NIỆM NỀN TẢNG
% ------------------------------------------------------------------------------
\begin{frame}{2. Phân biệt 3 Khái niệm Nền tảng (Tránh Đánh đồng)}
    \begin{table}[h]
        \centering
        \small
        \begin{tabularx}{\textwidth}{p{3.2cm} X X}
            \toprule
            \textbf{Khái niệm} & \textbf{Bản chất \& Góc nhìn} & \textbf{Ví dụ trong Quản lý Đào tạo} \\
            \midrule
            \textcolor{SoftTeal}{\textbf{1. User Needs}}\newline\textit{(Nhu cầu Người dùng)} & 
            Mong muốn chủ quan của người học nhằm hoàn thành lộ trình đào tạo. & 
            \textit{"Sinh viên muốn biết rõ lớp học còn chỉ tiêu và thời khóa biểu phù hợp để chọn ngay."} \\
            \midrule
            \textcolor{WarmAmber}{\textbf{2. User Problem}}\newline\textit{(Vấn đề Người dùng)} & 
            Trở ngại, rào cản thực tế ngăn cản người học đạt được nhu cầu. & 
            \textit{"Sinh viên bị nghẽn mạng, thiếu công cụ mô phỏng TKB, chọn nhầm môn trùng giờ."} \\
            \midrule
            \textcolor{ThemeNavy}{\textbf{3. Business Problem}}\newline\textit{(Vấn đề Nhà trường)} & 
            Tổn thất về nhân lực, kinh phí vận hành, lãng phí tài nguyên của tổ chức. & 
            \textit{"Phòng Đào tạo quá tải xử lý 1.500 đơn/kỳ; lãng phí 12\% công suất phòng học."} \\
            \bottomrule
        \end{tabularx}
    \end{table}
    
    \vspace{0.1cm}
    \begin{tcolorbox}[colback=CardBg, colframe=MutedSlate]
        \footnotesize \textbf{\textcolor{ThemeNavy}{Nguyên tắc then chốt:}} Vấn đề của người dùng dẫn đến hành vi tiêu cực (khiếu nại, giữ chỗ ảo), chính hành vi đó mới tạo thành tổn thất kinh doanh cho Nhà trường!
    \end{tcolorbox}
\end{frame}

% ------------------------------------------------------------------------------
% SLIDE 4: USER PROBLEM STATEMENT VỚI BẰNG CHỨNG
% ------------------------------------------------------------------------------
\begin{frame}{3. Câu phát biểu Vấn đề Người dùng (Dựa trên Bằng chứng)}
    \begin{tcolorbox}[colback=PrimaryBlue!5, colframe=PrimaryBlue, title=Cấu trúc Chuẩn hóa (User Problem Statement Formula)]
        \centering\small
        \textbf{[Persona]} cần \textbf{[Nhu cầu]} nhưng gặp trở ngại vì \textbf{[Rào cản]}, dẫn đến \textbf{[Hậu quả]}, được chứng minh qua \textbf{[Bằng chứng nghiên cứu]}.
    \end{tcolorbox}
    
    \vspace{0.15cm}
    
    \begin{tcolorbox}[colback=CardBg, colframe=ThemeNavy, title=Phát biểu Cụ thể cho Sinh viên]
        \small
        \textit{"\textbf{Sinh viên đại học theo học chế tín chỉ (Persona)} cần một công cụ trực quan để chủ động mô phỏng và kiểm tra tính tương thích của thời khóa biểu dự kiến \textbf{(Need)}, nhưng đang gặp trở ngại vì thông tin môn học phân tán và hạ tầng đăng ký thường xuyên quá tải vào giờ mở cổng \textbf{(Obstacle)}, dẫn đến việc bị trượt các môn học quan trọng và xáo trộn tiến độ tốt nghiệp \textbf{(Impact)}, được chứng minh qua \textbf{82\% sinh viên khảo sát từng gặp sự cố đăng ký mỗi đầu kỳ (Evidence)}."}
    \end{tcolorbox}
\end{frame}

% ------------------------------------------------------------------------------
% SLIDE 5: PROBLEM STATEMENT CANVAS (CHUẨN 6 Ô THEO BÀI GIẢNG)
% ------------------------------------------------------------------------------
\begin{frame}{4. Problem Statement Canvas (Khung Nghiên cứu 6 Khối Chuẩn)}
    \begin{columns}[T]
        % CỘT 1: CONTEXT & CUSTOMERS
        \begin{column}{0.32\textwidth}
            \begin{tcolorbox}[colback=CardBg, colframe=ThemeNavy, title=1. CONTEXT (Bối cảnh)]
                \scriptsize
                \textbf{Xảy ra khi nào?}\\
                • 2-3 tuần trước kỳ: Sinh viên tự lên kế hoạch học tập.\\
                • 8h00 sáng ngày mở cổng: Hàng chục ngàn kết nối cạnh tranh chỉ tiêu lớp học phần có hạn.
            \end{tcolorbox}
            
            \vspace{0.08cm}
            \begin{tcolorbox}[colback=CardBg, colframe=ThemeNavy, title=4. CUSTOMERS (Đối tượng)]
                \scriptsize
                \textbf{Ai chịu tác động nhiều nhất?}\\
                • Sinh viên năm 2, 3 (bước vào các môn chuyên ngành có môn học phụ thuộc).\\
                • Sinh viên cần học cải thiện điểm.\\
                • Sinh viên theo học chương trình kép.
            \end{tcolorbox}
        \end{column}

        % CỘT 2: PROBLEM & IMPACT
        \begin{column}{0.34\textwidth}
            \begin{tcolorbox}[colback=CardBg, colframe=WarmAmber, title=2. PROBLEM (Vấn đề \& Gốc rễ)]
                \scriptsize
                \textbf{Nguyên nhân cốt lõi?}\\
                • Thông tin môn học phân tán trên file tĩnh (PDF/Excel), thiếu tính tương tác.\\
                • Tiếp nhận đăng ký dồn cục; thiếu công cụ dự báo và kiểm tra trước khi xác nhận.
            \end{tcolorbox}
            
            \vspace{0.08cm}
            \begin{tcolorbox}[colback=CardBg, colframe=WarmAmber, title=5. IMPACT (Tác động)]
                \scriptsize
                \textbf{Emotional Impact:}\\
                • Căng thẳng, bất an, cảm giác "may rủi".\\
                \textbf{Quantifiable Impact:}\\
                • \textbf{82\%} sinh viên trượt môn dự kiến.\\
                • \textbf{1.5 - 3h} xếp hàng nộp đơn tại trường.\\
                • \textbf{1.500+} đơn khiếu nại/học kỳ.\\
                • \textbf{12\%} lớp ảo gây lãng phí phòng.
            \end{tcolorbox}
        \end{column}

        % CỘT 3: ALTERNATIVES & SHORTCOMINGS
        \begin{column}{0.32\textwidth}
            \begin{tcolorbox}[colback=CardBg, colframe=SoftTeal, title=3. ALTERNATIVES (Hiện tại)]
                \scriptsize
                \textbf{SV đang tự xoay xở thế nào?}\\
                • Tự kẻ bảng xếp lịch trên Excel/sổ tay.\\
                • Hỏi kinh nghiệm truyền miệng qua mạng.\\
                • Sử dụng tool/extension tự chế ngoài luồng.\\
                • Nộp đơn giấy xin mở lớp tại khoa.
            \end{tcolorbox}
            
            \vspace{0.08cm}
            \begin{tcolorbox}[colback=CardBg, colframe=SoftTeal, title=6. SHORTCOMINGS (Điểm yếu)]
                \scriptsize
                \textbf{Hạn chế của cách tự xoay xở?}\\
                • Excel: Tốn 4-6h, dễ nhầm kíp học.\\
                • Tool ngoài: Rò rỉ tài khoản, spam server.\\
                • Đơn giấy: Quá tải nhân sự, mất 2-3 tuần xử lý thủ công.
            \end{tcolorbox}
        \end{column}
    \end{columns}
\end{frame}

% ------------------------------------------------------------------------------
% SLIDE 6: KẾT NỐI NHÂN QUẢ (USER PROBLEM <-> BUSINESS PROBLEM)
% ------------------------------------------------------------------------------
\begin{frame}{5. Chuỗi Kết nối Nhân quả: Vấn đề Người dùng $\leftrightarrow$ Nhà trường}
    \begin{center}
        \begin{tikzpicture}[node distance=2.2cm, auto, >=stealth, thick]
            \node[draw=WarmAmber, fill=WarmAmber!10, rounded corners, text width=2.8cm, align=center] (up) 
                {\textbf{\small 1. Vấn đề Sinh viên}\\\scriptsize Trượt môn mong muốn, lịch học bị chồng chéo};
            
            \node[draw=PrimaryBlue, fill=PrimaryBlue!10, rounded corners, text width=2.8cm, align=center, right of=up, xshift=1.4cm] (act) 
                {\textbf{\small 2. Hành vi Nảy sinh}\\\scriptsize Nộp đơn giấy xin lớp, giữ chỗ lớp ảo rồi rút sau};
            
            \node[draw=SoftTeal, fill=SoftTeal!10, rounded corners, text width=2.8cm, align=center, right of=act, xshift=1.4cm] (imp) 
                {\textbf{\small 3. Tác động Vận hành}\\\scriptsize Bộ máy Một cửa quá tải; lớp biến động sĩ số};
            
            \node[draw=ThemeNavy, fill=ThemeNavy!10, rounded corners, text width=2.8cm, align=center, right of=imp, xshift=1.4cm] (bp) 
                {\textbf{\small 4. Vấn đề Nhà trường}\\\scriptsize Tốn ngân sách ngoài giờ; lãng phí giảng đường};
            
            \draw[->, line width=1.5pt, PrimaryBlue] (up) -- (act);
            \draw[->, line width=1.5pt, PrimaryBlue] (act) -- (imp);
            \draw[->, line width=1.5pt, PrimaryBlue] (imp) -- (bp);
        \end{tikzpicture}
    \end{center}
    
    \vspace{0.25cm}
    \begin{tcolorbox}[colback=CardBg, colframe=ThemeNavy]
        \footnotesize \textbf{Luận điểm Báo cáo:} Sinh viên không có ý định gây khó dễ cho Nhà trường. Chính rào cản trong việc lập kế hoạch đã buộc họ phải có những hành vi ứng phó tiêu cực. Giải quyết trải nghiệm lập kế hoạch cho sinh viên chính là cách triệt tiêu gốc rễ bài toán quá tải của Nhà trường!
    \end{tcolorbox}
\end{frame}

% ------------------------------------------------------------------------------
% SLIDE 7: LEAN UX CANVAS (PHẦN 1: Ô 1 - 5)
% ------------------------------------------------------------------------------
\begin{frame}{6. Lean UX Canvas (Phần 1: Khám phá Vấn đề \& Hướng Giải pháp)}
    \begin{columns}[T]
        \begin{column}{0.48\textwidth}
            \begin{tcolorbox}[colback=CardBg, colframe=ThemeNavy, title=Box 1: Vấn đề Tổ chức (Business Problem)]
                \footnotesize
                Mô hình tổ chức đào tạo bị động; tốn chi phí giải quyết hơn 1.500 đơn khiếu nại đầu kỳ; tỷ lệ lớp học phần ảo chiếm 12\% gây lãng phí phòng học.
            \end{tcolorbox}
            
            \vspace{0.08cm}
            \begin{tcolorbox}[colback=CardBg, colframe=SoftTeal, title=Box 2: Kết quả Kỳ vọng (Business Outcomes)]
                \footnotesize
                \begin{itemize}
                    \item Giảm \textbf{85\%} đơn khiếu nại và đề nghị điều chỉnh lớp.
                    \item Giảm tỷ lệ lớp học phần ảo từ \textbf{12\% xuống < 2\%}.
                    \item Ổn định thời khóa biểu toàn trường sau \textbf{3 ngày} mở cổng.
                \end{itemize}
            \end{tcolorbox}
        \end{column}
        
        \begin{column}{0.48\textwidth}
            \begin{tcolorbox}[colback=CardBg, colframe=PrimaryBlue, title=Box 3 \& 4: Users \& User Outcomes]
                \footnotesize
                \textbf{Đối tượng:} Sinh viên, Giảng viên, Cán bộ Đào tạo.\\
                \textbf{Lợi ích người dùng:}
                \begin{itemize}
                    \item Sinh viên an tâm lên kế hoạch, không còn "canh me".
                    \item Giảng viên nắm bắt sĩ số chính xác sớm để chuẩn bị bài.
                \end{itemize}
            \end{tcolorbox}
            
            \vspace{0.08cm}
            \begin{tcolorbox}[colback=CardBg, colframe=WarmAmber, title=Box 5: Hướng Giải pháp (Solutions)]
                \footnotesize
                \begin{itemize}
                    \item \textbf{Cổng Lập kế hoạch Học tập Chủ động:} Cho phép tạo giỏ môn dự kiến và tự động kiểm tra tương thích trước mở cổng.
                    \item \textbf{Dự báo Nhu cầu (Demand-driven):} Thu thập nguyện vọng học phần trước 4 tuần để mở đúng số lớp.
                \end{itemize}
            \end{tcolorbox}
        \end{column}
    \end{columns}
\end{frame}

% ------------------------------------------------------------------------------
% SLIDE 8: LEAN UX CANVAS (PHẦN 2: Ô 6 - 8)
% ------------------------------------------------------------------------------
\begin{frame}{7. Lean UX Canvas (Phần 2: Giả thuyết \& Kế hoạch Thử nghiệm)}
    \begin{tcolorbox}[colback=PrimaryBlue!5, colframe=PrimaryBlue, title=Box 6: Giả thuyết Cốt lõi (Core Hypothesis)]
        \small
        \textit{"Chúng tôi tin rằng \textbf{số lượng đơn khiếu nại điều chỉnh lớp sẽ giảm 80\%} nếu \textbf{Sinh viên} được \textbf{mô phỏng và kiểm tra trước tính hợp lệ của thời khóa biểu} thông qua \textbf{Cổng Lập kế hoạch Học tập Chủ động}."}
    \end{tcolorbox}
    
    \vspace{0.15cm}
    
    \begin{columns}[T]
        \begin{column}{0.48\textwidth}
            \begin{tcolorbox}[colback=CardBg, colframe=WarmAmber, title=Box 7: Rủi ro Lớn nhất Cần học Trước]
                \footnotesize
                \begin{itemize}
                    \item Liệu sinh viên có chủ động dành thời gian xếp lịch thử trước ngày mở cổng chính thức hay không?
                    \item Giao diện kiểm tra thời khóa biểu có đủ đơn giản để sinh viên không cần đọc hướng dẫn?
                \end{itemize}
            \end{tcolorbox}
        \end{column}
        
        \begin{column}{0.48\textwidth}
            \begin{tcolorbox}[colback=CardBg, colframe=SoftTeal, title=Box 8: Kế hoạch Thử nghiệm MVP Ít Tốn Công]
                \footnotesize
                \begin{enumerate}
                    \item \textbf{Bản mẫu Tương tác (Figma Prototype):} Xây dựng giao diện kéo thả môn học và cảnh báo lịch chồng chéo.
                    \item \textbf{Thử nghiệm Khả năng sử dụng (15 SV):} Đo thời gian sinh viên hoàn thành xếp 1 TKB mẫu.
                    \item \textbf{Khảo sát Mức độ hài lòng (SUS):} Đạt điểm $\ge 80/100$.
                \end{enumerate}
            \end{tcolorbox}
        \end{column}
    \end{columns}
\end{frame}

% ------------------------------------------------------------------------------
% SLIDE 9: TỔNG KẾT & BƯỚC TIẾP THEO
% ------------------------------------------------------------------------------
\begin{frame}{8. Đóng góp cho Dự án IT3120 \& Lời cảm ơn}
    \begin{tcolorbox}[colback=CardBg, colframe=ThemeNavy, title=Giá trị Đóng góp cho Đề tài Lớn IT3120]
        \small
        \begin{itemize}
            \item \textbf{Hoàn thiện vững chắc Chương 1:} Xác định bối cảnh nghiệp vụ theo chuỗi $S \to P \to I \to N \to V$ và chứng minh tính khả thi theo Khung DVF dựa trên bằng chứng khảo sát thực tế.
            \item \textbf{Định hướng chính xác cho Chương 2:} Từ nhu cầu người dùng đã kiểm chứng, nhóm tự tin mô hình hóa các Ca sử dụng (Use Cases) cốt lõi mà không sợ bị lạc đề hay dư thừa tính năng.
        \end{itemize}
    \end{tcolorbox}
    
    \vspace{0.35cm}
    \begin{center}
        \large \textbf{\textcolor{ThemeNavy}{XIN TRÂN TRỌNG CẢM ƠN THẦY CÔ VÀ CÁC BẠN!}}\\
        \vspace{0.15cm}
        \normalsize \textit{Mời Thầy Cô và các bạn đặt câu hỏi (Q\&A)}
    \end{center}
\end{frame}

\end{document}
